\section{Introduction}
\label{sec:intro}

Coding agents such as Claude Code now take part in everyday software work: they edit
files, run commands, and review changes on a developer's behalf. To make an agent follow a
team's procedure for a recurring task, developers package that procedure as an
\emph{agent skill}. A skill is a folder whose \texttt{SKILL.md} file starts with YAML front
matter, holding a \texttt{name} and a \texttt{description}, followed by instructions written
in Markdown; the folder may also contain scripts and reference documents~\cite{agentskills_spec,
destefanis2027gitskills}. The agent first reads only the descriptions of the available skills
and loads a skill's full instructions when it judges that the current task matches one.
Anthropic released the format as an open specification in October 2025, and nine months
later millions of \texttt{SKILL.md} files were sitting in public GitHub
repositories~\cite{destefanis2027gitskills}.

Skills spread differently from conventional dependencies. A library installed from npm or
PyPI is declared in a manifest, pinned to a version, and updated through the package manager.
A skill has no such channel. Developers copy its folder into their own repository, often edit
it along the way, and nothing records where it came from. The GitSkills dataset shows how
common this is: of 3,797,117 \texttt{SKILL.md} files collected from 282,200 repositories,
50.5\% are byte-for-byte copies of another file in the dataset~\cite{destefanis2027gitskills}.
Copies that were edited after copying are not part of that figure, because any edit produces
a new content hash.

We use the following terms throughout. An \emph{exact copy} has the same content hash as
another file. A \emph{derived variant} is a skill produced by copying another skill and then
editing it: renaming, trimming, extending, rewording, or translating it. Two skills copied
from a common source are derived variants of each other. A \emph{template sibling} shares only
generic scaffolding with another skill, such as section headings, boilerplate guidance, or a
catalog's house style, while its task-specific instructions were written separately. A
\emph{skill family} is a group of skills connected by derivation, and its \emph{lineage} is
the order in which its members appeared. These notions follow code-clone research, where
clones range from identical fragments (Type-1) to fragments with renamed identifiers
(Type-2) and with added, removed, or changed statements (Type-3)~\cite{roy2009comparison}.

Lineage matters because an agent executes a skill instead of merely displaying it. A person
reading an outdated README can notice the problem and work around it, whereas an agent acts
on stale instructions directly~\cite{gao2026registry}. Skills can also carry risk. Security
audits have found vulnerable and malicious skills in public collections~\cite{liu2026wild},
and in January 2026 a single campaign planted more than 1,200 malicious skills in the
OpenClaw skill marketplace~\cite{zhu2026skillreuse}. Once a flawed skill has been copied and
edited, its relatives are hard to find.
